% Main proof fragment for the full-normal, variable-source clause calculus.
% Requires amsmath, amssymb, amsthm and theorem/lemma/corollary environments.

\section{Equality needs only substitution for a numerical variable}
\label{sec:variable-equality}

The numerical terms are the inherited full computation normal forms of
\[
 t::=x\mid0\mid S(t)\mid I(n,a;u.h).
\]
Let $N$ denote the full normalizer. It contracts both computation rules
\[
 I(0,a;u.h)\longrightarrow a,\qquad
 I(Sn,a;u.h)\longrightarrow h[I(n,a;u.h)/u].
\]
All substitutions below avoid capture and are followed by $N$. We use
the inherited coherence of normalization with substitution, not an
unproved substitution theorem for derivations. Put $1=S0$, $2=SS0$,
and use the term abbreviation
\[
 C(v,a,b):=N(I(v,a;u.b)),\qquad u\notin\operatorname{FV}(b).
\]
Thus $C(0,a,b)\equiv a$ and $C(Sv,a,b)\equiv b$ are literal
canonical-syntax identities.

Formulas have grammar $A::=(s=t)\mid D_{x,y}A$. Contexts are finite
reusable lists of clauses $[A;y.B(y)]$, with $y$ bound only in the
consequent and absent from $A$. An ordinary assumption $A$ abbreviates
$[H;v.A]$, where $H=(0=0)$ and $v$ is unused. Its use follows from
reflexivity at zero and clause use.

\subsection{The complete rule inventory}

There are five inference patterns. The only changed one is VSub.
\[
 \frac{}{\Gamma\vdash t=t}.\tag{Refl}
\]
\[
 \frac{\Gamma\vdash x=t\qquad\Gamma\vdash u=v}
      {\Gamma\vdash v=N(u[t/x])},
 \qquad x\text{ is a free numerical variable}.\tag{VSub}
\]
\[
 \frac{\Gamma\vdash A}{\Gamma\vdash B(t)},
 \qquad[A;y.B(y)]\in\Gamma.\tag{Use}
\]
\[
 \frac{\Gamma\vdash D_{x,y}B(x,y)\qquad
       \Gamma,[B(b,0);y.B(b,y)]\vdash A}
      {\Gamma\vdash A}.\tag{E}
\]
\[
 \frac{\Gamma,B(t,\eta)\vdash B(t,Sz)}
      {\Gamma\vdash D_{x,y}B(x,y)},\qquad
 \eta=\begin{cases}
 z,&B\text{ is a literal equation and }t\equiv0,\\
 0,&\text{otherwise}.
 \end{cases}\tag{I}
\]
In E, $b$ is fresh for $\Gamma,A$ and the passive parameters of $B$.
In I, $z$ is fresh for $\Gamma,t$ and the passive matrix parameters.
Every displayed instance is canonical and well scoped. There are no
other logical, computation, equality, successor, or induction rules.

VSub substitutes simultaneously for every free occurrence of $x$ in
$u$. It neither deletes $x$ from the numerical context nor substitutes
in $\Gamma$ or $v$. The target $t$ may contain $x$: inserted targets
are not traversed again. No general term template, selected occurrence,
path, composite source, or equality-pattern matching is supplied.
The source can be any free numerical variable, not only the most
recently introduced one. Bound iterator state variables are not sources.

\subsection{Structural facts that do not require equality congruence}

Weakening, injective renaming of numerical parameters, and
ordinary-assumption cut are admissible by induction on derivations.
Variable sources remain variables under injective renaming, which
commutes with their capture-free substitution. The normalizer commutes
with renaming. Such renaming also preserves the introduction mode.
Before a weakening or cut, choose new eigenparameters away from the
finitely many inserted data. A cut replaces each use of the designated
constant clause by the supplied proof. These arguments do not assume
arbitrary substitution of numerical terms into derivations.

\subsection{Numerical cases before equality}

For a fresh unused coordinate $w$, abbreviate
\[
       q_yF(y):=D_{w,y}F(y).
\]

\begin{lemma}[Two conditional-closure interfaces]
The following rules are admissible using only I, E and Use:
\[
 \frac{\Gamma,F(0)\vdash F(Sz)}{\Gamma\vdash q_yF(y)},
 \qquad
 \frac{\Gamma\vdash q_yF(y)\qquad\Gamma\vdash F(0)}
      {\Gamma\vdash F(t)}.
\]
The first rule has the usual fresh $z$; the second permits every
well-scoped term $t$.
\end{lemma}
\begin{proof}
For introduction, use I at witness $1$ for the unused $w$ coordinate.
Because $1$ is not literally zero, I is in its case mode for every
matrix root. Its premise is exactly the displayed premise.

For elimination, E opens $[F(0);y.F(y)]$: its numerical witness is
unused in $F$. Use with the supplied base and target $t$ gives $F(t)$,
and E closes the branch. Freshening is allowed by the structural facts.
\end{proof}

\begin{corollary}[Numerical cases]
If $z$ is fresh for $\Gamma$ and the passive parameters of $F$, then
\[
 \frac{\Gamma\vdash F(0)\qquad\Gamma\vdash F(Sz)}
      {\Gamma\vdash F(t)}
\]
is admissible for every term $t$.
\end{corollary}
\begin{proof}
Weaken the successor proof by $F(0)$ and use closure introduction.
Specialize that closure with the given base. The target $t$ can contain
parameters already occurring in $\Gamma$; only the fresh generalized
parameter $z$ is excluded from $\Gamma$.
\end{proof}

No equality rule has been used. In particular, neither the repair of
the missing atomic case-introduction at witness zero nor a compound
envelope is needed for this numerical-case construction.

\subsection{A fresh name directly recovers every numerical template}

Define the parity term by
\[
 P(y):=N\bigl(I(y,0;u.C(u,1,0))\bigr).
\]
Only the two closed computations $P(1)\equiv1$ and $P(2)\equiv0$
will be used. No range theorem or induction for $P$ is required.

\begin{theorem}[Recovered unrestricted crossed transport]
For every numerical template $h(w)$ the rule
\[
 \frac{\Gamma\vdash s=t\qquad\Gamma\vdash N(h(s))=r}
      {\Gamma\vdash r=N(h(t))}
\]
is admissible. Its simulation uses exactly one instance of VSub.
\end{theorem}
\begin{proof}
Choose $x,y$ fresh for all passive data and form the equation matrix
\[
 B(x,y):=\bigl(
 C(y,0,C(P(y),x,N(h(x))))
 =C(y,0,C(P(y),t,r))\bigr).
\]
Introduce $D_{x,y}B(x,y)$ using witness $s$. Its successor goal is
\[
 C(P(Sz),s,N(h(s)))=C(P(Sz),t,r).
\]
Apply numerical cases to
\[
 F(v):=\bigl(C(v,s,N(h(s)))=C(v,t,r)\bigr).
\]
Its zero instance is literally the given $s=t$ and its successor
instance is literally the given $N(h(s))=r$. Numerical cases at
$P(Sz)$ therefore proves the required successor goal. Choose the
case eigenparameter fresh for the whole current local context and
independent of $z$. This proof ignores the temporary assumption
of I, so it works in both introduction modes, including when $s$
is literally zero.

Now use E to open the package at a fresh name $b$. The resulting
clause is $[H;y.B(b,y)]$, because $B(b,0)\equiv H$. Use that clause
at $1$ and at $2$, with reflexivity as the premise $H$. The two
closed parity computations give, literally in canonical syntax,
\[
 B(b,1)\equiv\bigl(N(h(b))=r\bigr),\qquad
 B(b,2)\equiv(b=t).
\]
One VSub, with major premise $b=t$ and minor premise $N(h(b))=r$,
yields
\[
 r=N\bigl(N(h(b))[t/b]\bigr)=N(h(t)),
\]
where the last displayed equality identifies canonical syntax by
normal-substitution coherence. It is not an invocation of an equality
rule. The fresh name $b$ occurs neither in the outer context nor in
the conclusion, so E closes.
\end{proof}

The entire argument preceding the single VSub uses only the logical
rules, reflexivity and literal computation. In particular, it assumes
no symmetry, transitivity, congruence, arbitrary numerical substitution
of derivations, equality reflection, or induction theorem. Passive
occurrences of $s$ in $h$ remain untouched, because $b$ was fresh for
$h$ and all other passive data. Normalization may duplicate or erase
the distinguished occurrences; the substitution-coherence identity
justifies this without counting occurrences.

As immediate derived rules, take $h(w)=w$ to obtain root joining,
use root joining and reflexivity for symmetry and transitivity, and
use crossed transport with reflexivity to obtain congruence. Two
crossed transports recover replacement in both sides of every
literal-equation template.

\subsection{Exact comparison and both interpretation composites}

Let $\mathcal C_T$ be the inherited fully computation-normal Eq/$D$
clause calculus with crossed transport, and let $\mathcal C_V$ be the
five-rule calculus above. Their canonical terms, formulas, clauses,
introduction modes and eigenvariable conditions are identical.

\begin{theorem}[Same complete canonical sequents]
For every canonical formula $A$ and finite canonical clause context
$\Gamma$,
\[
       \Gamma\vdash_{\mathcal C_V}A
       \quad\Longleftrightarrow\quad
       \Gamma\vdash_{\mathcal C_T}A.
\]
Both translations are effective on finite proof trees and are literally
the identity on the entire end sequent.
\end{theorem}
\begin{proof}
Forwards, VSub is one crossed transport. Given its source variable $x$,
abstract all free occurrences of $x$ in its minor left side $u$ by a
fresh template variable $w$, obtaining $K(w)$. Coherence gives
$N(K(x))=u$ and $N(K(t))=N(u[t/x])$. Crossed transport produces
exactly the VSub conclusion, including when $t$ itself contains $x$.
Translate all other nodes unchanged.

Backwards, replace each crossed-transport node by the preceding
derived construction after recursively translating its premises. Copy
all other rules. Each newly chosen name is fresh for the finite current
context, conclusion and term data; the fresh names used by the numerical-case construction are chosen
away from surrounding names. Hence all E and I
freshness conditions remain valid and the original sequent is unchanged.
\end{proof}

General numerical substitution on derivations follows by transfer
through this exact comparison. It is not a primitive-tag-preserving
operation on $\mathcal C_V$: a substitution may turn the variable
source of a VSub into a composite term. This is a derived admissibility
statement, proved through the recovered unrestricted transport.

Compose the exact comparison with the inherited normalized-clausal
equivalence and native/iPRA bi-interpretation. The ordinary realization
of $\mathcal C_V$ is therefore equiconsistent and bi-interpretable with
full first-order iPRA. All numerical domains are whole and equality
retains its actual meaning. Both numerical composite graphs remain
$x=y$; the formula round trips are those of the inherited comparisons.

\paragraph{What is reduced.}
The primitive equality source is now only a variable, and its action
is ordinary numerical substitution, which the binder calculus already
requires. The primitive equality inference accepts no additional term
template and performs no composite-pattern search. The required freedom
of ordinary equality is recovered by introducing, using and eliminating
fresh numerical names. This is stronger than merely excluding a literal
zero source. Full computation normalization is essential to the short
bootstrap presented here: it makes the zero, one and two instances of
the naming matrices literal identities before equality congruence is
available. The proof deliberately uses logical rules; it claims no
corresponding completeness theorem for the isolated equality fragment.


\subsection{The same reduction in the native six-rule presentation}

For completeness, the argument also applies before the clause
encoding. Keep canonical scalar terms, and use native formulas
\[
 A::=(s=t)\mid Q_yA\mid D_{x,y}A
\]
with ordinary formula contexts. The six patterns are assumption,
reflexivity, VSub, and the following three rules:
\[
 \frac{\Gamma\vdash Q_yB(y)\qquad\Gamma\vdash B(0)}
      {\Gamma\vdash B(t)};\tag{$Q$E}
\]
\[
 \frac{\Gamma\vdash D_{x,y}B(x,y)\qquad
       \Gamma,Q_yB(b,y)\vdash A}
      {\Gamma\vdash A};\tag{$D$E}
\]
\[
 \frac{\Gamma,B(t,\eta)\vdash B(t,Sz)}
      {\Gamma\vdash D_{x,y}B(x,y)}.\tag{$D$I}
\]
The mode $\eta$, freshness restrictions and VSub are exactly as
above. The $Q$ and $D$ formula constructors are primitive in this
presentation; no $Q$-introduction is postulated.

Numerical cases still uses no equality rule: introduce
$D_{x,y}F(y)$ at witness $1$ with unused $x$, ignoring the temporary
$F(0)$ assumption; its $D$E branch supplies $Q_yF(y)$, and $Q$E
with the given $F(0)$ supplies $F(t)$. For the naming matrix $B$,
$D$I and $D$E work exactly as in the theorem. The two uses of the
clause $[H;y.B(b,y)]$ are replaced by two $Q$E applications to the
ordinary assumption $Q_yB(b,y)$ and reflexivity $H$. The final single
VSub is unchanged. Conversely native VSub is a native crossed
transport as above. Thus the six-rule native VSub and crossed-transport
presentations prove exactly the same sequents under the identity
translation. This avoids any dependence on the clause encoding in
the measured native implementation.

\paragraph{Why the source is an arbitrary free variable.}
Restricting VSub to the newest numerical parameter would make direct
weakening of an arbitrary existing proof fail: inserting a fresh
parameter turns its former newest source into an older variable.
The proof above does not assume such a stronger restriction. Allowing
any free variable makes the required injective renamings and weakening
immediate, while still removing arbitrary composite sources and
supplied equality templates.
